% Options for packages loaded elsewhere
% Options for packages loaded elsewhere
\PassOptionsToPackage{unicode}{hyperref}
\PassOptionsToPackage{hyphens}{url}
\PassOptionsToPackage{dvipsnames,svgnames,x11names}{xcolor}
%
\documentclass[
  letterpaper,
  DIV=11,
  numbers=noendperiod]{scrartcl}
\usepackage{xcolor}
\usepackage{amsmath,amssymb}
\setcounter{secnumdepth}{-\maxdimen} % remove section numbering
\usepackage{iftex}
\ifPDFTeX
  \usepackage[T1]{fontenc}
  \usepackage[utf8]{inputenc}
  \usepackage{textcomp} % provide euro and other symbols
\else % if luatex or xetex
  \usepackage{unicode-math} % this also loads fontspec
  \defaultfontfeatures{Scale=MatchLowercase}
  \defaultfontfeatures[\rmfamily]{Ligatures=TeX,Scale=1}
\fi
\usepackage{lmodern}
\ifPDFTeX\else
  % xetex/luatex font selection
\fi
% Use upquote if available, for straight quotes in verbatim environments
\IfFileExists{upquote.sty}{\usepackage{upquote}}{}
\IfFileExists{microtype.sty}{% use microtype if available
  \usepackage[]{microtype}
  \UseMicrotypeSet[protrusion]{basicmath} % disable protrusion for tt fonts
}{}
\makeatletter
\@ifundefined{KOMAClassName}{% if non-KOMA class
  \IfFileExists{parskip.sty}{%
    \usepackage{parskip}
  }{% else
    \setlength{\parindent}{0pt}
    \setlength{\parskip}{6pt plus 2pt minus 1pt}}
}{% if KOMA class
  \KOMAoptions{parskip=half}}
\makeatother
% Make \paragraph and \subparagraph free-standing
\makeatletter
\ifx\paragraph\undefined\else
  \let\oldparagraph\paragraph
  \renewcommand{\paragraph}{
    \@ifstar
      \xxxParagraphStar
      \xxxParagraphNoStar
  }
  \newcommand{\xxxParagraphStar}[1]{\oldparagraph*{#1}\mbox{}}
  \newcommand{\xxxParagraphNoStar}[1]{\oldparagraph{#1}\mbox{}}
\fi
\ifx\subparagraph\undefined\else
  \let\oldsubparagraph\subparagraph
  \renewcommand{\subparagraph}{
    \@ifstar
      \xxxSubParagraphStar
      \xxxSubParagraphNoStar
  }
  \newcommand{\xxxSubParagraphStar}[1]{\oldsubparagraph*{#1}\mbox{}}
  \newcommand{\xxxSubParagraphNoStar}[1]{\oldsubparagraph{#1}\mbox{}}
\fi
\makeatother


\usepackage{longtable,booktabs,array}
\usepackage{calc} % for calculating minipage widths
% Correct order of tables after \paragraph or \subparagraph
\usepackage{etoolbox}
\makeatletter
\patchcmd\longtable{\par}{\if@noskipsec\mbox{}\fi\par}{}{}
\makeatother
% Allow footnotes in longtable head/foot
\IfFileExists{footnotehyper.sty}{\usepackage{footnotehyper}}{\usepackage{footnote}}
\makesavenoteenv{longtable}
\usepackage{graphicx}
\makeatletter
\newsavebox\pandoc@box
\newcommand*\pandocbounded[1]{% scales image to fit in text height/width
  \sbox\pandoc@box{#1}%
  \Gscale@div\@tempa{\textheight}{\dimexpr\ht\pandoc@box+\dp\pandoc@box\relax}%
  \Gscale@div\@tempb{\linewidth}{\wd\pandoc@box}%
  \ifdim\@tempb\p@<\@tempa\p@\let\@tempa\@tempb\fi% select the smaller of both
  \ifdim\@tempa\p@<\p@\scalebox{\@tempa}{\usebox\pandoc@box}%
  \else\usebox{\pandoc@box}%
  \fi%
}
% Set default figure placement to htbp
\def\fps@figure{htbp}
\makeatother


% definitions for citeproc citations
\NewDocumentCommand\citeproctext{}{}
\NewDocumentCommand\citeproc{mm}{%
  \begingroup\def\citeproctext{#2}\cite{#1}\endgroup}
\makeatletter
 % allow citations to break across lines
 \let\@cite@ofmt\@firstofone
 % avoid brackets around text for \cite:
 \def\@biblabel#1{}
 \def\@cite#1#2{{#1\if@tempswa , #2\fi}}
\makeatother
\newlength{\cslhangindent}
\setlength{\cslhangindent}{1.5em}
\newlength{\csllabelwidth}
\setlength{\csllabelwidth}{3em}
\newenvironment{CSLReferences}[2] % #1 hanging-indent, #2 entry-spacing
 {\begin{list}{}{%
  \setlength{\itemindent}{0pt}
  \setlength{\leftmargin}{0pt}
  \setlength{\parsep}{0pt}
  % turn on hanging indent if param 1 is 1
  \ifodd #1
   \setlength{\leftmargin}{\cslhangindent}
   \setlength{\itemindent}{-1\cslhangindent}
  \fi
  % set entry spacing
  \setlength{\itemsep}{#2\baselineskip}}}
 {\end{list}}
\usepackage{calc}
\newcommand{\CSLBlock}[1]{\hfill\break\parbox[t]{\linewidth}{\strut\ignorespaces#1\strut}}
\newcommand{\CSLLeftMargin}[1]{\parbox[t]{\csllabelwidth}{\strut#1\strut}}
\newcommand{\CSLRightInline}[1]{\parbox[t]{\linewidth - \csllabelwidth}{\strut#1\strut}}
\newcommand{\CSLIndent}[1]{\hspace{\cslhangindent}#1}



\setlength{\emergencystretch}{3em} % prevent overfull lines

\providecommand{\tightlist}{%
  \setlength{\itemsep}{0pt}\setlength{\parskip}{0pt}}



 


\KOMAoption{captions}{tableheading}
\makeatletter
\@ifpackageloaded{caption}{}{\usepackage{caption}}
\AtBeginDocument{%
\ifdefined\contentsname
  \renewcommand*\contentsname{Table of contents}
\else
  \newcommand\contentsname{Table of contents}
\fi
\ifdefined\listfigurename
  \renewcommand*\listfigurename{List of Figures}
\else
  \newcommand\listfigurename{List of Figures}
\fi
\ifdefined\listtablename
  \renewcommand*\listtablename{List of Tables}
\else
  \newcommand\listtablename{List of Tables}
\fi
\ifdefined\figurename
  \renewcommand*\figurename{Figure}
\else
  \newcommand\figurename{Figure}
\fi
\ifdefined\tablename
  \renewcommand*\tablename{Table}
\else
  \newcommand\tablename{Table}
\fi
}
\@ifpackageloaded{float}{}{\usepackage{float}}
\floatstyle{ruled}
\@ifundefined{c@chapter}{\newfloat{codelisting}{h}{lop}}{\newfloat{codelisting}{h}{lop}[chapter]}
\floatname{codelisting}{Listing}
\newcommand*\listoflistings{\listof{codelisting}{List of Listings}}
\makeatother
\makeatletter
\makeatother
\makeatletter
\@ifpackageloaded{caption}{}{\usepackage{caption}}
\@ifpackageloaded{subcaption}{}{\usepackage{subcaption}}
\makeatother
\usepackage{bookmark}
\IfFileExists{xurl.sty}{\usepackage{xurl}}{} % add URL line breaks if available
\urlstyle{same}
\hypersetup{
  pdftitle={Multidimensional Attribution of Coral Bleaching Mortality on the Great Barrier Reef},
  pdfauthor={Samuel Matthews},
  colorlinks=true,
  linkcolor={blue},
  filecolor={Maroon},
  citecolor={Blue},
  urlcolor={Blue},
  pdfcreator={LaTeX via pandoc}}


\title{Multidimensional Attribution of Coral Bleaching Mortality on the
Great Barrier Reef}
\author{Samuel Matthews}
\date{}
\begin{document}
\maketitle
\begin{abstract}
{[}Abstract Placeholder{]}
\end{abstract}


\section{Introduction}\label{introduction}

Climate change is increasing the frequency and intensity of marine
heatwaves, causing widespread coral bleaching and mortality across the
Great Barrier Reef (GBR) (Eakin, Sweatman, and Brainard 2019). As
thermal stress shifts from acute events to chronic disturbance, managers
must transition from monitoring to targeted spatial interventions (Oppen
et al. 2017). Deploying resilience-based strategies---like localized
solar radiation management, assisted gene flow, or culling benthic
competitors---requires accurate attribution of coral mortality to
specific thermal thresholds (Hughes et al. 2018; Oppen et al. 2017).
Treating the GBR uniformly masks the spatial variation in bleaching and
survival observed among neighboring reefs (Wold et al. 2023). Directing
conservation resources at actionable scales requires spatially explicit
models that identify the environmental and biological factors moderating
survival during severe heatwaves.

Current diagnostic tools rely on static or coarse-scale environmental
proxies that miss fine-scale physical and ecological mediators of
thermal stress (Safaie et al. 2018). Many vulnerability assessments
evaluate Degree Heating Weeks (DHW) as an isolated, linear driver of
mortality without adjusting for background coral recovery (MacNeil et
al. 2019). This omission conflates baseline demographic turnover with
bleaching-attributable loss, biasing mortality estimates. Most models
also omit localized optical and hydrodynamic modifiers---such as water
clarity (Secchi depth), cloud cover shading, and current flushing
speeds---or shifts in community composition, such as the proportion of
sensitive \emph{Acropora} taxa (Cacciapaglia and Woesik 2016).
Deterministic forecasts trained only on historical heatwaves often fail
under novel thermal regimes. This inflexibility led to widespread
underprediction of mortality during the 2024 mass bleaching event,
limiting the effectiveness of adaptive management.

We developed a growth-adjusted framework to quantify coral mortality
across the GBR. We combined multidecadal benthic photo-transects and
manta tow data with high-resolution environmental predictors, including
winter temperature variability, historical thermal exposure, and acute
optical modifiers. To isolate bleaching impacts, we adjusted baseline
mortality estimates using reef-specific Gompertz coral growth parameters
to account for expected biological recovery between surveys. We paired
machine-learning algorithms (Random Forests and Boosted Regression
Trees) to capture non-linear feature interactions with Bayesian
hierarchical beta regression models to bound proportional mortality. The
Bayesian models propagate parameter uncertainty and partition spatial
variance across reef clusters. By explicitly modeling interactions
between thermal stress and optical mediators---specifically how
localized Secchi depth and cloud fraction buffer DHW---we link mortality
attribution to measurable physical contexts.

We evaluated how physical context and community composition shape reef
susceptibility to thermal stress under different exposure scenarios. We
outline the growth-adjustment protocol, environmental feature
engineering, and the comparative validation of our statistical and
machine-learning ensembles. Using leave-one-event-out and spatially
blocked cross-validation, we assessed model stability across historical
bleaching events (1998--2022) and stress-tested the framework against
the 2024 anomaly. Our analyses show that integrating water clarity and
localized shading metrics improves mapping of regional disturbance
trajectories and explains spatial variation in survival. Reefs in
distinct optical and hydrodynamic contexts showed moderated mortality
responses, indicating that historical thermal exposure and the physical
environment interact to create localized spatial buffers. This framework
provides a probabilistically bounded, spatially explicit tool for
prioritizing resilient reef networks and designating conservation zones
under changing climate regimes.

\section{Methods}\label{methods}

\subsection{Benthic Data and
Growth-Adjustment}\label{benthic-data-and-growth-adjustment}

We integrated multidecadal benthic photo-transect and broad-scale manta
tow data across the Great Barrier Reef spanning 1998 to 2024. This
expanded dataset (\(N=465\) reef-event pairs) captures the full spectrum
of recent mass bleaching events. To isolate thermal impacts from
background demographic turnover, we growth-adjusted baseline mortality
estimates using reef-specific Gompertz coral growth parameters. This
adjustment accounts for the biological recovery expected between survey
intervals. We also constructed a strictly filtered subset of confirmed
bleaching mortality (\(N=305\)) to evaluate model sensitivity when
excluding compound disturbances such as cyclones and freshwater plumes.

\subsection{Environmental Covariates and Feature
Engineering}\label{environmental-covariates-and-feature-engineering}

We engineered a suite of high-resolution environmental and biological
predictors. Thermal stress was parameterized using Degree Heating Weeks
(DHW). To test for physical buffering, we integrated localized optical
and hydrodynamic modifiers, specifically acute water clarity (Secchi
depth), cloud cover fraction, and prevailing current speeds. We also
included winter sea surface temperature (SST) variability and the
preceding-year proportion of thermally sensitive \emph{Acropora} taxa to
account for community composition.

\subsection{Statistical Modeling and
Validation}\label{statistical-modeling-and-validation}

We applied a dual analytical architecture. We used Boosted Regression
Trees (BRT) to map non-linear feature interactions and derive relative
variable importance. We paired this with Bayesian hierarchical Beta and
Binomial-OLRE regression models to estimate dose-response trajectories
and enforce physical bounds on proportional mortality. The Bayesian
framework propagates parameter uncertainty and partitions spatial
variance across nested reef clusters. Model stability was assessed using
leave-one-event-out and spatial sector-blocked cross-validation, and
predictive skill was stress-tested on out-of-sample data from the
extreme 2024 mass bleaching event.

\begin{verbatim}
Warning: package 'readr' was built under R version 4.5.1
\end{verbatim}

\begin{verbatim}
Warning: package 'knitr' was built under R version 4.5.3
\end{verbatim}

\section{Results}\label{results}

\subsection{Model Validation and Variable
Importance}\label{model-validation-and-variable-importance}

The dual analytical framework reliably predicted proportional coral
mortality across historical bleaching eras. Spatial sector-blocked
cross-validation showed the Bayesian models achieved a strong predictive
\(R^2\) and minimized mean absolute error compared to isolated
univariate frameworks. Boosted Regression Trees identified Degree
Heating Weeks as the primary driver of mortality. However, community
composition (proportion \emph{Acropora}) and optical modifiers (Secchi
depth and cloud cover) contributed substantial relative influence,
confirming that thermal stress does not act in isolation.

\subsection{Dose-Response Trajectories and Environmental
Buffering}\label{dose-response-trajectories-and-environmental-buffering}

Bayesian dose-response curves demonstrated that physical context
strongly moderates reef susceptibility to thermal anomalies. High cloud
cover and reduced water clarity (lower Secchi depth) attenuated
mortality at severe DHW thresholds. This confirms that localized shading
and turbidity buffer acute thermal exposure. Conversely, reefs dominated
by \emph{Acropora} in clear-water, unshaded environments experienced
steeper mortality trajectories, matching the known physiological
sensitivity of these taxa.

\subsection{2024 Mass Bleaching
Generalization}\label{mass-bleaching-generalization}

When stress-tested against the unprecedented 2024 mass bleaching event
as a strict out-of-sample forecast, the framework maintained predictive
accuracy. While traditional DHW-only linear models severely
underpredicted 2024 mortality, integrating optical modifiers and
community composition allowed the BRT and Bayesian models to map the
extreme spatial heterogeneity observed during the anomaly.

\begin{longtable}[]{@{}llrrr@{}}

\caption{\label{tbl-summary}Summary of model performance metrics across
validation folds.}

\tabularnewline

\toprule\noalign{}
Dataset & Model & RMSE & Predictive R2 & Bias \\
\midrule\noalign{}
\endhead
\bottomrule\noalign{}
\endlastfoot
ltmp & benchmark & 0.130 & 0.397 & -0.004 \\
ltmp & brms & 0.133 & 0.360 & 0.010 \\
ltmp & brt & 0.140 & 0.297 & -0.017 \\
ltmp & formal\_ensemble & 0.131 & 0.388 & 0.000 \\
ltmp & stacked\_ensemble & 0.127 & 0.421 & -0.001 \\
manta & benchmark & 0.159 & -0.066 & -0.017 \\
manta & brms & 0.141 & 0.169 & -0.011 \\
manta & brt & 0.181 & -0.376 & -0.013 \\
manta & formal\_ensemble & 0.141 & 0.169 & -0.011 \\
manta & stacked\_ensemble & 0.141 & 0.169 & -0.011 \\

\end{longtable}

\section{Discussion}\label{discussion}

\section{References}\label{references}

\phantomsection\label{refs}
\begin{CSLReferences}{1}{0}
\bibitem[\citeproctext]{ref-cacciapaglia2016}
Cacciapaglia, Chris, and Robert van Woesik. 2016. {``Climate-Change
Refugia: Shading Reef Corals by Turbidity.''} \emph{Global Change
Biology} 22 (3): 1145--54. \url{https://doi.org/10.1111/gcb.13166}.

\bibitem[\citeproctext]{ref-eakin2019}
Eakin, C Mark, Hugh PA Sweatman, and Russell E Brainard. 2019. {``The
2014--2017 Global-Scale Coral Bleaching Event: Insights and Impacts.''}
\emph{Coral Reefs} 38 (4): 539--45.
\url{https://doi.org/10.1007/s00338-019-01844-2}.

\bibitem[\citeproctext]{ref-hughes2018}
Hughes, Terry P, James T Kerry, Andrew H Baird, Sean R Connolly, Andreas
Dietzel, C Mark Eakin, Scott F Heron, et al. 2018. {``Spatial and
Temporal Patterns of Mass Bleaching of Corals in the Anthropocene.''}
\emph{Science} 359 (6371): 80--83.
\url{https://doi.org/10.1126/science.aan8048}.

\bibitem[\citeproctext]{ref-macneil2019}
MacNeil, M Aaron, Camille Mellin, Samuel Matthews, Nicholas H Wolff, Tim
R McClanahan, Michelle Devlin, Christopher Drovandi, Kerrie Mengersen,
and Russ C Babcock. 2019. {``Water Quality Mediates Resilience on the
Great Barrier Reef.''} \emph{Nature Ecology \& Evolution} 3 (4):
620--27. \url{https://doi.org/10.1038/s41559-019-0832-3}.

\bibitem[\citeproctext]{ref-vanoppen2017}
Oppen, Madeleine JH van, Ruth D Gates, Linda L Blackall, Neal Cantin,
Leela J Chakravarti, Wing Y Chan, Craig Cormick, et al. 2017.
{``Shifting Paradigms in Restoration of the World's Coral Reefs.''}
\emph{Global Change Biology} 23 (9): 3437--48.
\url{https://doi.org/10.1111/gcb.13647}.

\bibitem[\citeproctext]{ref-safaie2018}
Safaie, Aryan, Nyssa J Silbiger, Tim R McClanahan, Geno Pawlak, Daniel J
Barshis, James L Hench, Justin S Rogers, Gareth J Williams, and Kristen
A Davis. 2018. {``High-Frequency Temperature Variability Reduces the
Risk of Coral Bleaching.''} \emph{Nature Communications} 9 (1): 1671.
\url{https://doi.org/10.1038/s41467-018-04074-2}.

\bibitem[\citeproctext]{ref-wold2023}
Wold, Renee, Jacquelyn Simpson, Spencer Miller, Joshua R Hancock,
Crawford Drury, and Joshua S Madin. 2023. {``Fine-Scale Variability in
Coral Bleaching and Mortality During a Marine Heatwave.''}
\emph{Frontiers in Marine Science} 10: 1108365.
\url{https://doi.org/10.3389/fmars.2023.1108365}.

\end{CSLReferences}




\end{document}
